%%
%% fall-lecture23.tex
%% 
%% Made by Alex Nelson <pqnelson@gmail.com>
%% Login   <alex@lisp>
%% 
%% Started on  2026-11-19T08:40:09-0800
%% Last update 2026-11-19T08:40:09-0800
%% 

\lecture[Some inequalities]{}

We will finish last lecture's topics (especially important is Fatou's
lemma), then start the next topic ``inequalities''.

\begin{lemma}[Fatou]
Let $f_{n}\colon X\to[0,\infty]$ be a sequence of non-negative
measurable functions on $(X,\mathcal{A},\mu)$. Then
\begin{equation}
\int(\liminf f_{n})\leq\liminf\int f_{n}.
\end{equation}
\end{lemma}

\begin{proof}
We see that we can construct the increasing sequence of measurable functions
\begin{equation}
g_{n} := \inf_{k\geq n}f_{k}
\end{equation}
Then
\begin{subequations}
  \begin{align}
\int\liminf_{n}f_{n} &= \int\lim_{n\to\infty}\inf_{k\geq n}f_{k}\\
&=\int\lim_{n}g_{n}\quad\mbox{by def of $g_{n}$}\\
&=\lim_{n}\int g_{n}\quad\mbox{by monotone convergence theorem}.
  \end{align}
\end{subequations}
We also see
\begin{subequations}
  \begin{align}
    \lim_{n}\int g_{n} &= \liminf_{n}\int g_{n}\\
    &= \liminf_{n}\int(\inf_{k\geq n}f_{k})\\
    &\leq \liminf_{n}\int f_{n},
  \end{align}
\end{subequations}
since $g_{n}\leq f_{n}$ by definition of infinimum.
\end{proof}

\begin{example}
To remember this, take $f_{n}=\mathbf{1}_{[n,\infty)]}$ be the
indicator function on $[n,\infty)$ for each $n\in\NN$. Then
  $\int\liminf f_{n}=0$ and $\liminf\int f_{n}=\infty$ and $0\leq\infty$.
\end{example}

\begin{theorem}[Lebesgue Dominated Convergence]
Let $f_{n}\xrightarrow{a.e.}f$ be a sequence of measurable functions.
Let $g\in L^{1}(\mu)$. If $|f_{n}|\leq g$ for all $n$, then
$f_{n}\xrightarrow{L^{1}}f$ in norm.
\end{theorem}

\begin{proof}
Well, $\int g<\infty$ is finite, so
\begin{equation}
\int 2g=\int\lim_{n}(2g-|f_{n}-f|)
\end{equation}
since
\begin{equation}
\lim_{n}(2g-|f_{n}-f|)\stackrel{a.e.}{=}2g,
\end{equation}
since $f_{n}\xrightarrow{a.e.}f$ by hypothesis.
But now,
\begin{equation}
2g-|f_{n}-f|\geq0,
\end{equation}
since
\begin{equation}
|f_{n}-f|\leq|f_{n}|\leq g\leq 2g\quad\mbox{a.e.,}
\end{equation}
by domination.

Now, we can apply Fatou's Lemma. [The factor of 2 in $2g$ is because
we have $f_{n}$ and $f$]
Then we see
\begin{subequations}
  \begin{align}
\int\lim_{n}(2g-|f_{n}-f|) &= \int\liminf_{n}(2g-|f_{n}-f|)\\
&\leq\liminf\int(2g-|f_{n}-f|)\quad\mbox{by Fatou},
  \end{align}
  but
  \begin{align}
\int2g &= \int\lim_{n}(2g-|f_{n}-f|) \\
&\leq \liminf_{n}\int(2g-|f_{n}-f|)\\
&\leq(\int2g)-\limsup_{n}\int|f_{n}-f|,
  \end{align}
\end{subequations}
which means that
\begin{equation}
\limsup_{n}\int|f_{n}-f|\leq0,
\end{equation}
so $\limsup\|f_{n}-f\|_{L^{1}}\leq0$. But since the norm is non-negative,
$\limsup\|f_{n}-f\|_{L^{1}}=0$ which implies
\begin{equation}
\lim_{n\to\infty}\|f_{n}-f\|_{L^{1}}=0.
\end{equation}
Hence the theorem.
\end{proof}

\begin{fact}[Chebyshev's inequality]
Let $f\in L^{1}(\mu)$. Then for any $\varepsilon>0$ we have
\begin{equation}
\mu\{|f|>\varepsilon\}\leq\frac{\|f\|_{L^{1}}}{\varepsilon}.
\end{equation}
\end{fact}

\begin{proposition}
If $f_{n}\xrightarrow{L^{1}}f$ converges in $L^{1}$-norm,
then $f_{n}\xrightarrow{\mu}f$ converges in measure.
\end{proposition}

\begin{proof}
Let $E_{\varepsilon}=\{|f|>\varepsilon\}$. Then
\begin{subequations}
  \begin{align}
\|f\|_{L^{1}} &= \int_{X}|f|\,\D\mu\\
&\geq\int_{E_{\varepsilon}}|f|\,\D\mu\\
&\geq\int_{E_{\varepsilon}}\varepsilon\,\D\mu=\varepsilon\mu(E_{\varepsilon}),
  \end{align}
\end{subequations}
which is precisely the claim after dividing through by $\varepsilon$.
\end{proof}

\subsection{Inequalities}

\begin{node}
Estimates and inequalities lie at the heart of the applications of
$L^{p}$ spaces in analysis. The most basic of these are H\"{o}lder and
Minkowski inequalities. We will start with a warm-up: Jensen's inequality.
First, we need to review/define the notion of convex functions.
\end{node}

\begin{definition}
A function $\varphi\colon(a,b)\to\RR$ is \define{Convex} if for all
$\lambda\in(0,1)$ we have $\varphi(\lambda x_{1}+(1-\lambda)x_{2})\leq\lambda\varphi(x_{1})+(1-\lambda)\varphi(x_{2})$
for any $a<x_{1}<x_{2}<b$.
\end{definition}

\begin{node}
Equivalently, $\varphi\colon(a,b)\to\RR$ is convex if the region
$E=\{(x,y)\mid y\geq\varphi(x)\}$ is a convex region, meaning any two
points of $E$ can be connected by a straight line entirely contained
in $E$. If we define
\begin{equation}
\operatorname{slope}_{x_{1},x_{2}}(\varphi):=\frac{\varphi(x_{2})-\varphi(x_{1})}{x_{2}-x_{1}}
\end{equation}
for the slope of the line segment connecting $(x_{1},\varphi(x_{1}))$
to $(x_{2},\varphi(x_{2}))$, then for any $a<x_{1}<x_{2}<x_{3}<b$ we have
\begin{equation}
\operatorname{slope}_{x_{1},x_{2}}(\varphi)\leq
\operatorname{slope}_{x_{1},x_{3}}(\varphi)\leq
\operatorname{slope}_{x_{2},x_{3}}(\varphi).
\end{equation}
\end{node}

\begin{fact}
Convex functions are continuous (and therefore Borel measurable).
\end{fact}

\begin{fact}
Twice differentiable functions $f\in C^{2}(a,b)$ are convex if their
second derivatives are non-negative $f''\geq0$ on its domain.
\end{fact}

\begin{theorem}[Jensen's inequality]
Let $\mu$ be a probability measure on set $\Omega$ [so $\mu(\Omega)=1$]
and $f\colon\Omega\to(a,b)$ be integrable $f\in L^{1}(\mu)$.
Let $\varphi\colon(a,b)\to\RR$ be convex.
Then $\varphi(\int f\,\D\mu)\leq\int_{\Omega}(\varphi\circ f)\,\D\mu$.
\end{theorem}

[Fact: $\varphi$ is continuous, so $\varphi\circ f$ is measurable.]

\begin{proof}
Denote $t=\int f\,\D\mu$ for the expected value of $f$, and clearly $a<t<b$.
Also denote,
\begin{equation}
\beta:=\sup_{s\in(a,b)}\frac{\varphi(t)-\varphi(s)}{s-t}.
\end{equation}
We have the inequality for any $t<u<b$,
\begin{subequations}
\begin{equation}
\beta\leq\frac{\varphi(u)-\varphi(t)}{u-t},
\end{equation}
and for any $a<s<t$,
\begin{equation}
\frac{\varphi(t)-\varphi(s)}{s-t}\leq\beta.
\end{equation}
\end{subequations}
Combining these statements, we see for any $s\in(a,b)$,
\begin{equation}
\varphi(s)\geq\varphi(t)-\beta(t-s).
\end{equation}
Now, we set $s=f(x)$, and then integrate against $\D\mu(x)$ in this inequality
\begin{subequations}
  \begin{align}
\int\varphi(f(x))\,\D\mu &\geq \int(\varphi(t)-\beta(t-f(x)))\,\D\mu\\
&=(\int\varphi(t))-\beta(t-\int f)\\
&=(\int\varphi(t))-\beta(t-t)\\
&=\int\varphi(t)\\
&=\varphi(t)\mu(\Omega)=\varphi(t)\cdot1=\varphi(t).
  \end{align}
\end{subequations}
Hence the claim.
\end{proof}